The constitutional authority for federal redistricting legislation is well established, though its precise scope has been contested.

\subsection{Article I, \S4: The Elections Clause}

Article I, \S4, cl.\ 1 provides: ``The Times, Places and Manner of holding Elections for Senators and Representatives, shall be prescribed in each State by the Legislature thereof; but the Congress may at any time by Law make or alter such Regulations, except as to the Place of chusing Senators.'' This clause is the primary constitutional basis for the DIA.

The Elections Clause's breadth is not in doubt. In \textit{Ex parte Siebold}, 100 U.S.\ 371 (1880), the Supreme Court held that Congress may supplement, displace, or correct state election regulations. In \textit{Smiley v. Holm}, 285 U.S.\ 355 (1932), the Court confirmed that Congress could constitutionally mandate that states draw congressional districts as single-member geographic units. In \textit{Cook v. Gralike}, 531 U.S.\ 510 (2001), the Court clarified that the clause grants states authority to regulate ``the procedural mechanisms for obtaining a ballot position,'' and grants Congress equivalent authority to override state rules.

The Elections Clause question most relevant to the DIA is whether Congress may specify a particular method---not merely criteria---for how states draw their lines. The apportionment analogy (Section~5) suggests yes: 2 U.S.C.\ \S2a specifies the Huntington-Hill algorithm with mathematical precision, and it has never been constitutionally challenged. The DIA applies the same structural template to within-state districting.

\subsection{Section 5 of the Fourteenth Amendment}

The DIA also asserts authority under Section 5 of the Fourteenth Amendment, which grants Congress power to enforce the Amendment's requirements by appropriate legislation. \textit{Reynolds v. Sims}, 377 U.S.\ 533 (1964), established that the Equal Protection Clause requires substantially equal population among congressional districts. Congress may enforce this requirement through legislation.

The dual-authority design in DIA \S101(f) is a defense against the risk that a future Court narrows the Elections Clause. If either authority is sufficient---Article I or Fourteenth Amendment \S5---the statute survives even if one basis is rejected. This design follows the pattern of the Voting Rights Act, which also asserts dual authority under the Fifteenth Amendment and the Fourteenth.

\subsection{The Printz/Murphy Anti-Commandeering Limit}

The most significant constitutional risk for the DIA is the anti-commandeering doctrine. \textit{Printz v. United States}, 521 U.S.\ 898 (1997), holds that Congress may not ``commandeer'' state officers to administer federal programs. \textit{Murphy v. NCAA}, 138 S.\ Ct.\ 1461 (2018), extended this principle to prohibit Congress from ``dictating'' that states maintain a particular regulatory regime.

The DIA's rebuttal is textual: the Elections Clause \emph{expressly} authorizes federal alteration of state election regulations. As DIA \S101(g) explains, the Tenth Amendment does not constrain powers the Constitution affirmatively assigns to Congress. This argument is supported by \textit{Foster v. Love}, 522 U.S.\ 67 (1997), which upheld a federal statute overriding Louisiana's open-primary system, and by \textit{Branch v. Smith}, 538 U.S.\ 254 (2003), which allowed federal courts to draw congressional maps when states failed to timely comply with federal requirements.

The anti-commandeering risk can be further reduced by structuring the DIA as a direct federal mandate (addressed to states as sovereign entities acting in their governmental capacity) rather than as a directive to individual state officers. This distinction matters under \textit{Printz}'s framework.

\subsection{Non-Delegation Concerns}

The DIA delegates to the Director of the Census Bureau the authority to specify certain technical parameters---the seed for the algorithm, the version of the reference binary, and the TIGER adjacency errata process. This delegation is narrowly circumscribed: the statute itself specifies the algorithm (recursive bisection), the population metric (decennial census counts), the balance tolerance (\textpm 0.5\%), and the edge-weight specification. The Director's authority extends only to release-specific administrative decisions that cannot be fixed in statute.

Under \textit{West Virginia v. EPA}, 597 U.S.\ 697 (2022)'s major questions doctrine, regulations of significant political importance require clear congressional authorization. The DIA provides that authorization explicitly, by name, with specific parameter ranges. The delegation is ministerial, not policymaking---analogous to the Census Director's existing authority to specify the format of decennial data releases.
